% DAFx26 Demo Track Paper: FreeEQ8/ProEQ8
% Completely standalone - no custom .sty or .bst files needed

\documentclass[10pt,twocolumn,a4paper]{article}

% Standard packages only
\usepackage[utf8]{inputenc}
\usepackage[T1]{fontenc}
\usepackage{times}
\usepackage{amsmath,amssymb}
\usepackage{booktabs}
\usepackage{url}
\usepackage{hyperref}
\usepackage[margin=20mm,top=25mm,bottom=25mm]{geometry}

% Paper metadata
\title{\textbf{REAL-TIME STATE-SPACE PARAMETERIZATION IN DIGITAL EQUALIZATION: A PRODUCTION-GRADE SVF IMPLEMENTATION}}

\author{Gary Doman\\
Independent Researcher (GareBear99 / TizWildin)\\
\url{https://github.com/GareBear99/FreeEQ8}\\
\texttt{neovectr.inc@gmail.com}}

\date{}

\begin{document}
\maketitle

\begin{abstract}
We present FreeEQ8/ProEQ8, an open-source parametric equalizer using the Simper State Variable Filter (SVF) topology for modulation-stable Dynamic EQ. Both SVF and RBJ use BLT with identical prewarping and produce identical steady-state responses; SVF chosen for lower automation noise, better coefficient interpolation, improved SNR near DC. The architecture features lock-free SPSC triple-buffering for real-time safety, a variable-cadence coefficient engine reducing Dynamic EQ CPU cost by 75\%, and allocation-free semantic resonance detection. Measured performance: 0.62\% single-core CPU at 44.1\,kHz (161$\times$ headroom). Source code, benchmarks, and CI validation available under GPL-3.0.
\end{abstract}

\section{Introduction}
Traditional parametric EQ implementations use the Robert Bristow-Johnson (RBJ) Audio EQ Cookbook biquad formulas [1] in Transposed Direct Form II. While mathematically correct, the bilinear transform introduces frequency cramping near Nyquist: a Bell filter at 16\,kHz with $Q=1.0$ at 44.1\,kHz exhibits an effective bandwidth \textbf{199\% narrower} than intended.

Industry solutions include brute-force oversampling (adding latency and CPU cost) or proprietary polynomial analog-matching curves. This paper documents a third path: the Simper SVF topology [2], which pre-warps the cutoff frequency via $g = \tan(\pi \cdot f_c / f_s)$, achieving exact cutoff placement without oversampling.

\subsection{Product Architecture}
The codebase produces two plugins from a single source tree:

\textbf{FreeEQ8} (Free, GPL-3.0): 8-band parametric EQ using RBJ TDF-II biquads. Zero audio restrictions during real-time playback.

\textbf{ProEQ8} (\$20): 24-band parametric EQ using Simper SVF with modulation-stable SVF topology. Adds 4 saturation modes, A/B comparison, and collision detection.

\section{Filter Topology}

\subsection{RBJ TDF-II (FreeEQ8)}
Standard biquad transfer function with coefficients per the RBJ cookbook. 64-bit double internal state, float I/O.

\begin{table}[ht]
  \caption{RBJ Q distortion vs frequency (Bell, $Q=1.0$, 44.1\,kHz).}
  \centering
  \begin{tabular}{ccc}\toprule
    Frequency & Effective Q & Error \\\midrule
    1\,kHz & 1.005 & +0.5\% \\
    8\,kHz & 1.337 & +33.7\% \\
    16\,kHz & 2.990 & +199\% \\\bottomrule
  \end{tabular}
  \label{tab:qdistortion}
\end{table}

\subsection{Simper SVF (ProEQ8)}
Pre-warped cutoff via trapezoidal integration [2]:
\begin{equation}
g = \tan(\pi \cdot f_c / f_s), \quad k = 1/Q
\end{equation}
\begin{equation}
a_1 = 1 / (1 + g(g + k)), \quad a_2 = g \cdot a_1, \quad a_3 = g \cdot a_2
\end{equation}

\begin{table}[ht]
  \caption{SVF vs RBJ magnitude response at 16\,kHz.}
  \centering
  \begin{tabular}{ccc}\toprule
    Topology & Magnitude at $f_c$ & Error \\\midrule
    RBJ @ 44.1\,kHz & +6.000\,dB & $0.000$\,dB (identical to SVF) \\
    SVF @ 44.1\,kHz & +6.00\,dB & 0.00\,dB \\
    RBJ @ 4$\times$ OS & +5.993\,dB & $-0.007$\,dB \\\bottomrule
  \end{tabular}
  \label{tab:svfvsrbj}
\end{table}

\section{Real-Time Safety Architecture}

\subsection{SPSC Triple-Buffer}
Three buffer slots indexed by \{writeSlot, midSlot, readSlot\}. Audio thread writes and atomically swaps writeSlot $\leftrightarrow$ midSlot with \texttt{memory\_order\_release}. UI thread reads via midSlot $\leftrightarrow$ readSlot with \texttt{memory\_order\_acquire}. No mutex, no blocking.

\subsection{Variable-Cadence Dynamic EQ}
Dynamic EQ recomputes coefficients per-sample when active. During sustained signals, envelope changes $< 0.1$\,dB are inaudible within a 4-sample window (0.09\,ms). \textbf{Result:} 75--80\% reduction in coefficient updates.

\section{Benchmarks}
All benchmarks from standalone C++ (no JUCE, no DAW). Platform: x86-64, g++ -O3.

\begin{table}[ht]
  \caption{Performance (44.1\,kHz, 512-sample blocks).}
  \centering
  \begin{tabular}{ccc}\toprule
    Configuration & ns/sample & CPU Headroom \\\midrule
    RBJ 8-band stereo & 40.7 & 277$\times$ \\
    SVF 8-band stereo & 70.4 & 161$\times$ \\
    SVF DynEQ worst-case & 370.9 & 30.6$\times$ \\\bottomrule
  \end{tabular}
\end{table}

Concurrent stress tests (400\,ms, Intel i7-3720QM): 239M samples, \textbf{0 data tears}.

\section{Demonstration Plan}
The live demonstration will showcase: (1) Side-by-side RBJ vs SVF frequency response at 16\,kHz; (2) Real-time spectrum analyzer with lock-free triple-buffer rendering; (3) Dynamic EQ envelope tracking on drum transients; (4) 24-band surgical EQ workflow in ProEQ8.

\textbf{Requirements:} Table with power outlet, external audio interface, headphones.

\section{Conclusions}
The SVF topology is used for modulation-stable Dynamic EQ. Source code and CI validation (pluginval strictness-level-10) available at \url{https://github.com/GareBear99/FreeEQ8} under GPL-3.0.

\section{Acknowledgments}
The Simper SVF topology is due to Andrew Simper (Cytomic). The RBJ biquad coefficients and Audio EQ Cookbook are due to Robert Bristow-Johnson.

The 5-DOF constraint framework, the geometric-mean pinning approach, and the $Q\cdot\sqrt{G}$ convention are drawn from public commentary by Robert Bristow-Johnson on r/DSP (May 2026) in response to questions raised by this project. His analysis identified that cramping is a BLT property affecting both topologies equally.

The characterization of the SVF's actual advantages --- improved SNR near DC, reduced coefficient-change noise, and stable interpolation for Dynamic EQ --- is due to SkoomaDentist (r/DSP, May 2026). JUCE forum feedback from Nitsuj70, holy-city, zsliu98, and kerfuffle identified errors in the original paper. This work is released under GPL-3.0.

\section*{References}
\begin{enumerate}
\item R. Bristow-Johnson, ``Audio EQ Cookbook,'' musicdsp.org, 1994.
\item A. Simper, ``Solving the continuous SVF equations using trapezoidal integration,'' Cytomic, 2013.
\item W. Pirkle, \textit{Designing Audio Effect Plugins in C++}, 2nd ed., Focal Press, 2019.
\item R. Bristow-Johnson, comment on r/DSP ``I wrote a paper on why your EQ plugin lies to you,'' Reddit, May 2026. \url{https://www.reddit.com/r/DSP/comments/1tqynr5/}
\item S.~P.~Orfanidis, ``Digital parametric equalizer design with prescribed Nyquist-frequency gain,'' \textit{J. Audio Eng. Soc.}, vol.~45, no.~6, 1997.
\item R. Bristow-Johnson, ``Why does the Direct Form I become our first choice for fixed-point implementation?'' DSP Stack Exchange, 2022. \url{https://dsp.stackexchange.com/a/76855}
\item R. Bristow-Johnson, ``Best implementation of a real-time, fixed-point IIR filter,'' DSP Stack Exchange. \url{https://dsp.stackexchange.com/a/21811}
\end{enumerate}

\end{document}
